\documentclass[10pt]{article}

\usepackage[utf8]{inputenc}

\usepackage[T1]{fontenc}

\usepackage{graphicx}

\usepackage[export]{adjustbox}

\graphicspath{ {./images/} }

\usepackage{amsmath}

\usepackage{amsfonts}

\usepackage{amssymb}

\usepackage[version=4]{mhchem}

\usepackage{stmaryrd}



\begin{document}

ные дополнительными структурами (напр., отношением включения бинарных отношений, включения или равенства областей определения, включения или равенства образов и т. д.).



Лuт.: [1] Л я пи н Е. С., Полугруппы, М., 1960; [2] К лиффор д А. Х., П р е ст о н Г.Б., Алгёбраическая



\includegraphics[max width=\textwidth, center]{2023_07_08_6bdc953d30cd4afd1278g-1}

S c h e i n B. M., "Semigroup Forum", $1970, \mathrm{v}, 1, N_{0} 1, \mathrm{p} .1-62$.



Л. М. Глускин, Е. С. Ллпин. гомотопич. топологии; инвариант, равный нулю, если соответствующая задача разрешима, и отличный от нуля - в противном случае.



Пусть $(X, A)$ - пара клеточных пространств и $Y-$ односвязное (более общо - гомотопически простое) топологич. пространство. Можно ли данное непрерывное отображение $g: A \rightarrow Y$ продолжить до непрерывного отображения $f: X \rightarrow Y$ ? Продолжение можно осуществлять по остовам $X^{n}$ пространства $X$. Пусть построено такое отображение $f: X^{n} \cup A \rightarrow Y$, что $\left.f\right|_{A}=g$. Для любой ориентированной $(n+1)$-мерной клетки $S^{n} \rightarrow Y$ отображение $\left.f\right|_{\partial e^{n+1}}$ задает отоб ражение $S^{n} \rightarrow Y$ (где $S^{n}$ есть $n$-мерная сфера) и элемент $\alpha_{e} \in \pi_{n}(Y)$ (именно здесь и используется гомотопич. простота пространства $\boldsymbol{Y}$, позволяющая игнорирова ть выбор отмеченной точки). Таким образом, возни кает коцепь



\[

c_{f}^{n+1} \in C^{n+1}\left(X ; \pi_{n}(Y)\right), c_{f}^{n+1}\left(e^{n+1}\right)=\alpha_{e} .

\]



Так как для $e^{n+1} \subset A$, очевидно, $c_{f}^{n+1}\left(e^{n+1}\right)=0$, то на самом деле



\[

c_{f}^{n+1} \in C^{n+1}\left(X, A ; \pi_{n}(Y)\right) .

\]



Очевидно, $. c_{f}^{n+1}=0$ тогда и только тогда, когда $f$ продолжается на $X^{n+1}$, т. е. коцепь $c_{f}^{n+1}$ является п р еп я т с т в и е м к продолжению $f$ на $X^{n+1}$.



Коцепь $c_{f}^{n+1} \in C^{n+1}\left(X, A ; \pi_{n}(Y)\right)$ является коциклом. Из того, что $c_{f}^{n+1}=0$, вообще говоря, не следует, что $g$ не продолжается на $X: f$ может не продолжаться на $X^{n+1}$ в силу веудачного выбора продолжения $g$ на $X^{n}$. Может оказаться, напр., что отображение $f \mid X^{n-1} \cup A$ продолжается на $X^{n+1}$, т. е. что продолжение возможно после отступления на один шаг. Оказывается, что П. к этому является класс когомологий



\[

\left[c_{f}^{n+1}\right] \in H^{n+1}\left(X, A ; \pi_{n}(Y)\right),

\]



то есть $\left[c_{f}^{n+1}\right]=0$ тогда и только тогда, когда существует такое отображение $\tilde{f}: X^{n+1} \cup A \rightarrow Y$, что $\tilde{f} \mid X^{n-1} \cup A=$ $=f X_{X^{n-1} \cup A}$ (в частности, $\left.\tilde{f}\right|_{A}=g$ ). Для доказательства әтого утверждения используется конструкция различающей.



Так как задачу гомотопич. классификации отображений $X \rightarrow Y$ можно интерпретировать как задачу продолжения, то теория П. применима и к описанию множества $[X, Y]$ гомотопич. классов отображений из $X$ в $Y$. Пусть $I=[0,1]$ и пусть $A=X \times\{0,1\}$ - подпространство в $X \times I$. Тогда пара отображений $f_{0}, f_{1}: X \rightarrow Y$ интерпретируется как отображение $G: A \rightarrow Y, G(x, i)=$ $=f_{i}(x), i=0,1$, и наличие гомотопии между $f_{0}$ и $f_{1}$ означает наличие отображения $F: X \times I \rightarrow Y$, продолжающего отображение $G$. При этом если гомотопия $F$ построена на $n$-мерном остове пространства $X$, то П. к ее продолжению на $X$ есть различающая



\[

d^{n}\left(f_{0}, f_{1}\right) \in C^{n}\left(X ; \pi_{n}(Y)\right) .

\]



В качестве приложения можно указать описание множества $[X, Y]=-[X, K(\pi, n)], n>1$, где $K(\pi, n)-\ni \check{u}-$ ленберга - Маклейна пространство: $\pi_{i}(K(\pi, n))=0$ при $i \neq n, \pi_{n}(K(\pi, n))=\pi$. Пусть $f_{0}: X \rightarrow K(\pi, n)-$ постоянное отображение, а $f: X \rightarrow K(\pi, n)$ - произ- вольное непрерывное отображение. Так как $H^{i}(X$; $\left.\pi_{i}(Y)\right)=0$ при $i<n$, то $f_{0}$ и $f_{1}$ гомотопны на $X^{n-1}$ и можно, выбрав какую-нибудь такую гомотопию, определить различающую



\[

d^{n}\left(f, f_{0}\right) \in C^{n}\left(X ; \pi_{n}(Y)\right)=C^{n}(X ; \pi) .

\]



Класс когомологий $\left[d^{n}\left(f, f_{0}\right)\right] \in H^{n}(X ; \pi)$ определен корректно, т. е. не зависит от выбора гомотопии между $f_{0}$ и $f$ (в силу того, что $\pi_{i}(Y)=0$ при $i<n$ ). Далее, если отображения $f, g: X \rightarrow Y$ таковы, что $\left[d^{n}\left(f, f_{0}\right)\right]=$ $=\left[d^{n}\left(g, f_{0}\right)\right]$, то $\left[d^{n}(f, g)\right]=0$ и, значит, $f$ и $g$ гомотопны на $X^{n}$. П. к продолжению этой гомотопии на все $X$ лежат в группах $H^{i}\left(X ; \pi_{i}(Y)\right)=0$ (так как $\left.i>n\right)$, и, значит, отображения $f$ и $g$ гомотопны. Таким образом, гомотопич. класс отображения $f$ полностью определяется элементом $\left[d^{n}\left(f, f_{0}\right)\right] \in H^{n}(X, \pi)$. Наконец, для любого әлемента $x \in H^{n}(X, \pi)$ найдется отображение $f \quad$ с $\left[d^{n}\left(f, f_{0}\right)\right]=x$ и потому $[X, K(\pi, n)]=H^{n}(X ; \pi)$. Аналогично: если $\pi_{i}(Y)=i$ при $i<n$ и $\operatorname{dim} X \leqslant n$, то $[X, Y]=$ $=H^{n}\left(X ; \pi_{n}(Y)\right)$.



IІ ри исследовании задачи продолжения рассматривалась возможность продолжения после «отступления на один шаг». Полное решение задачи требует анализа возможности продолжения после отступления на произвольное число ппагов. Для этой цели используются когомологические операции или Постникова системы. Так, для описания множества $[X, Y]$, где $\pi_{i}(Y)=0$ при $i<n, \pi_{n}(Y) \neq 0, \operatorname{dim} X=n+r ;$ требуется, вообце говоря, исследовать возможность отступления на $r+1$ шаг, для чего надо исследовать первые $n+r$ этажей системы Постникова пространства $Y$, т. е. использовать когомологич. операции порядков $\leqslant r$ (в ст. Когомологические операции эта задача разобрана для $r=1$ ).



Теория П. используется также в более общей ситуации продолжения сечений. Пусть $p: E \rightarrow B$ - нек-рое расслоение со слоем $F$ (причем $\pi_{1}(F)=0$ и $\pi_{1}(B)$ действует на $\pi_{1}(F)$ тривиально), пусть $A \subset B$ и $s: A \rightarrow E-$ нек-рое сечение (т. е. такое непрерывное отображение, что $p s(a)=a)$. Можно ли продолжить $s$ на все $B$ ? Соответствующие П. лежат в группах $H^{n+1}\left(B ; \pi_{n}(F)\right)$. Задача продолжения получается из этой задачи, если положить $B=X, \quad E=X \times Y, p(x, y)=x, s(a)=(a, g(a))$. Аналогичным образом с помощью теории П. исследуется и задача классификации сечений.



Наконец, в задаче продолжения можно снять ограничение гомотопич. простоты пространства $Y$ (и аналогично в задаче о сечения $)$; для этого надо использовать когомологии с локальными коэффициентами.



ПРЕСЛЕДОВАНИЯ ИГРА - антагонистическая дифференциальная игра преследователя (догоняющего) $P$ и преследуемого (убегаюпего) $E$, движения к-рых описываются системами дифференциальных уравнений:



\[

P: \dot{x}=f(x, u), E: \dot{y}=g(y, v) \text {, }

\]



где $x, y$ - фазовые векторы, определяющие состояния игроков $P$ и $E$ соответственно; $u, v-$ управляющие параметры, выбираемые игроками в каждый момент времени из заданных компактных множеств $U, V$ евклидовых пространств. Целью $P$ может быть, напр., сближение с $E$ на заданное расстояние, что формально означает попадание $x$ в $l$-окрестность $y(l \geqslant 0)$. При этом различаются случаи сближения за минимальное время (П. и. н а 6 ы с т р о д е й с т в и е), к заданному моменту времени (П. И. с п р е д п и с а н н о й п р о д о л ж т е л ь н о с т ю ю) и до момента достижения игроком $E$ нек-рого множества (игра с «линией жизни»). Сравнительно хорошо изучены игры с полной информацией, когда оба игрока знают фазовые состояния друг друга в каждый текущий момент времени. Под решением П. и. понимается нахождение ситуации равновесия.





\end{document}